% Compile with: pdflatex neural_network.tex
% This is a complete document. To embed just the diagram in an article,
% copy its preamble definitions and tikzpicture into your article instead.
\documentclass[tikz]{article}
\usepackage{amsmath,amssymb}
\usepackage{listofitems} % Indexed lists, including two-dimensional matrices.
\usepackage{xcolor}
\usepackage{url}
\usepackage{tikz}
\usetikzlibrary{arrows.meta,backgrounds}
\usepackage{geometry}
\geometry{margin=1.5in}
% \geometry{landscape}

% ================================================================
% NETWORK DATA: edit these values to change the diagram.
% ================================================================
% List positions are 1-based; mathematical layers are 0-based:
% position \lay=1 means input layer \ell=0, etc.
\readlist\Nnod{4,6,4,3,2}

% EXAMPLE PARAMETERS ONLY: edit the weights, biases, and activations below.
% W^(ell) has one ROW per destination neuron i and one COLUMN per
% source neuron j: z^(ell) = W^(ell) a^(ell-1) + b^(ell).
% Therefore the edge a_j^(ell-1) -> a_i^(ell) uses w_ij^(ell).
% Each non-input neuron also has one bias; each non-input layer has one
% activation function, applied elementwise to its pre-activation vector.
%
% \setsepchar{;/,}: semicolons separate rows; commas separate columns.
% Example: a 2-by-3 matrix is entered as {1,2,3;4,5,6}.
\setsepchar{;/,}

% \csname W1\endcsname is a control sequence whose name is W1.
% Ordinary \W1 would instead mean the command \W followed by digit 1.
% \expandafter expands \csname...\endcsname before \readlist sees it.
\newcommand{\DefineWeightMatrix}[2]{%
    \expandafter\readlist\csname W#1\endcsname{#2}%
}

% W^(1): 6 destination neurons by 4 source neurons.
\DefineWeightMatrix{1}{%
     1, 0,-1, 2;
     0, 1, 2,-1;
    -1, 2, 1, 0;
     2,-1, 0, 1;
     1, 1,-2, 0;
     0,-2, 1, 2%
}
% W^(2): 4 destination neurons by 6 source neurons.
\DefineWeightMatrix{2}{%
     1, 0,-1, 2, 1, 0;
     0, 1, 2,-1, 0, 1;
    -1, 2, 0, 1,-2, 1;
     2,-1, 1, 0, 1,-2%
}
% W^(3): 3 destination neurons by 4 source neurons.
\DefineWeightMatrix{3}{%
     1,-1, 0, 2;
     0, 2, 1,-1;
    -2, 0, 1, 1%
}
% W^(4): 2 destination neurons by 3 source neurons.
\DefineWeightMatrix{4}{%
     1,-1, 2;
    -2, 1, 0%
}

% BIASES: a list of lists, one inner list for each NON-INPUT layer.
% Outer index = mathematical layer ell (1 through 4).
% Inner index = destination neuron i (1 through Nnod[ell+1]).
% Semicolons separate layers; commas separate neurons within a layer.
% Thus b^(1) has 6 entries, b^(2) has 4, b^(3) has 3, b^(4) has 2.
% For example, \Bias{1}{3} returns -1, the third bias in layer 1.
\readlist\Biases{%
     0, 1,-1, 0, 2,-2;
     1, 0,-1, 2;
     0, 1,-1;
     1, 0%
}
\setsepchar{,} % Restore the ordinary separator for later lists.

% ACTIVATIONS: a flat list, one function for each NON-INPUT layer.
% No dummy input-layer entry is needed: list index = ell, starting at 1.
% Keep compound expressions inside braces if they contain a comma.
% These are math-mode function names for typesetting, not a numerical
% forward-pass evaluator. Changing this list changes the layer headings.
\newcommand{\ReLU}{\operatorname{ReLU}}
\readlist\ActivationFunctions{\ReLU,\ReLU,\tanh,\sigma}

% Dynamic lookup: \Weight{ell}{i}{j} returns matrix entry W^(ell)[i,j].
% For example, \Weight{1}{2}{3} = -1 (row 2, column 3 of W^(1)).
% \csname builds the appropriate matrix name from the layer number.
\newcommand{\Weight}[3]{\csname W#1\endcsname[#2,#3]}
% Both accessors below also use mathematical, not list-position, layers.
\newcommand{\Bias}[2]{\Biases[#1,#2]}
\newcommand{\Activation}[1]{\ActivationFunctions[#1]}

% ================================================================
% COLORS: every active color below is black or white.
% To restore the original palette, replace the right-hand side of
% each \colorlet with the expression given in its inline comment.
% Then restore the style options indicated beside the styles below.
% ================================================================
\colorlet{ColorOutput}{black}       % Restore: red!80!black
\colorlet{ColorHidden}{black}      % Restore: blue!80!black
\colorlet{ColorInput}{black}     % Restore: green!60!black
\colorlet{myorange}{black}    % Restore: orange!70!red!60!black
\colorlet{ColorOutputDark}{black}   % Restore: red!30!black
\colorlet{ColorHiddenDark}{black}  % Restore: blue!40!black
\colorlet{ColorInputDark}{black} % Restore: green!30!black

% ================================================================
% EXPLANATION 1: WHAT DO THE STYLES MEAN?
% ================================================================
% \tikzset{...} defines reusable sets of TikZ drawing options.
% name/.style={...} means "when I use [name], apply these options."
% A style can include another style: "node in" starts with "node".
%
% thick             : use TikZ's thick stroke width (normally 0.8pt).
% circle            : give the neuron a circular outline.
% draw=black        : stroke its boundary in black.
% text=black        : typeset its contents in black.
% fill=white        : paint its interior white.
% minimum size=26pt : minimum node width and height, not its radius.
% inner sep=0.5pt   : padding between the text box and shape boundary.
% outer sep=0.6pt   : extra clearance used when connecting to the node.
% Explicit units (pt) replace the original implicit point dimensions.
%
% Original color mixes such as ColorInput!25 mean 25% ColorInput + 75% white.
% A bare color option such as blue!20!black sets the drawing/text color;
% draw=..., text=..., and fill=... let you control them separately.
\tikzset{
    % EXPLANATION 2: >=latex sets the default END arrow-tip shorthand >.
    % It does not draw an arrow by itself; -> would use that default.
    % Modern equivalent used here: >=Latex from arrows.meta.
    % The explicit -{Latex[...]} below overrides this default per path.
    >=Latex,
    node/.style={
        thick,circle,draw=black,text=black,fill=white,
        minimum size=26pt,inner sep=0.5pt,outer sep=0.6pt
    }, % Restore the original base outline with draw=ColorHidden.
    node in/.style={node,text=black,draw=black,fill=white},
    % Restore above: text=green!20!black,draw=ColorInput!30!black,fill=ColorInput!25
    node hidden/.style={node,text=black,draw=black,fill=white},
    % Restore above: text=blue!20!black,draw=ColorHidden!30!black,fill=ColorHidden!20
    node convol/.style={node,text=black,draw=black,fill=white},
    % Restore above: text=orange!20!black,draw=myorange!30!black,fill=myorange!20
    % "node convol" is retained from your file but is unused here:
    % this diagram is fully connected and does not depict convolution.
    node out/.style={node,text=black,draw=black,fill=white},
    % Restore above: text=red!20!black,draw=ColorOutput!30!black,fill=ColorOutput!20
    connect/.style={thick,draw=black}, % Restore: draw=ColorHiddenDark
    connect arrow/.style={
        connect,-{Latex[length=4pt,width=3.5pt]},
        shorten <=0.5pt,shorten >=1pt
    },
    % -{Latex[...]}: arrowhead only at the destination, with given size.
    % shorten <=... retracts the START; shorten >=... retracts the END.
    % Optional line cap=round rounds exposed stroke ends; it is cosmetic.
    node 1/.style={node in},     % Style 1 = input neurons.
    node 2/.style={node hidden}, % Style 2 = hidden neurons.
    node 3/.style={node out},    % Style 3 = output neurons.
    layer label/.style={align=center,text=black,font=\small},
    bias label/.style={text=black,fill=white,inner sep=1pt,font=\scriptsize},
    % Optional bias color: add text=ColorHiddenDark; keep fill=white for clarity.
    formula/.style={align=center,text=black,font=\small},
    % Optional formula color: replace text=black with text=ColorHiddenDark.
    weight label/.style={text=black,fill=white,inner sep=0.6pt,font=\scriptsize}
    % Restore weight-text color above with text=ColorHiddenDark; keep fill=white.
}

% EXPLANATION 3: WHAT DOES THIS EXPRESSION DO?
% \def stores a PGF math expression; it is evaluated inside the loop.
% Syntax: condition ? value_if_true : value_if_false.
% If \lay < \Nnodlen (not the final layer), use min(2,\lay):
%   \lay=1 -> 1 (input style); \lay=2,3,4 -> 2 (hidden style).
% Otherwise use 3 (output style). int(...) makes the result an integer,
% so that "node \n" becomes "node 1", "node 2", or "node 3".
% \Nnodlen is provided by listofitems and is 5 for this architecture.
\def\nstyle{int(\lay<\Nnodlen ? min(2,\lay) : 3)}

\begin{document}

\section{AI Generated}
GPT 6.1 Sol High Effort

\url{https://chatgpt.com/g/g-p-6aab0741ce3c8191b7ba47a512862502-machine-learning/c/6aedb7fd-3ea8-83e8-9d12-58bc3f930671}


\section{Prompt}
Use the attached latex as the basis for a neural network diagram but make the following changes:

Label every layer with $\ell$ like I started attempting

Define and dynamically lookup the weights that should be used instead of 100.

Explain the comments I had enumerated I didn't understand.

Change all of the colors to regular black and white text, but annotate what would need to change in order to restore the colors with inline comments

\begin{tikzpicture}[x=3cm,y=1.5cm]
    % x and y specify the physical lengths of one coordinate unit.
    % The wider layer spacing accommodates all 66 numerical edge labels.

    % PASS 1: create all neurons and one heading per layer.
    \foreachitem \N \in \Nnod{
        % \N is the current list value; \Ncnt is its 1-based position.
        % \edef expands \Ncnt immediately, freezing that position in \lay.
        \edef\lay{\Ncnt}
        \pgfmathtruncatemacro{\layer}{\lay-1}
        \pgfmathtruncatemacro{\n}{\nstyle}

        \foreach \i in {1,...,\N}{
            % Center EVERY column at y=0, regardless of its node count.
            \pgfmathsetmacro{\y}{(\N+1)/2-\i}
            % (N\lay-\i) names this node for later references.
            \node[node \n] (N\lay-\i) at (\layer,\y)
                {$a_{\i}^{(\layer)}$};
            \ifnum\lay>1
                % One bias per destination neuron, dynamically looked up.
                % anchor=north attaches the TOP of this small text box to
                % a point 2pt below the neuron's south (bottom) anchor.
                % The bias is added before the activation; see the formula.
                \node[bias label,anchor=north,yshift=-2pt]
                    at (N\lay-\i.south)
                    {$b_{\i}^{(\layer)}=\Bias{\layer}{\i}$};
            \fi
        }

        % This label is OUTSIDE the neuron loop: one label per layer.
        % All headings share the same y-coordinate for alignment.
        % \ell is the mathematical symbol; \else is the TeX branch keyword.
        % The original \ellse was a typo and would not compile.
        \ifnum\lay=1
            \node[layer label] at (\layer,3.6)
                {Input layer\\$\ell=\layer$\\$\mathbf{a}^{(0)}=\mathbf{x}$};
            % Restore color: add text=ColorInput!60!black to this node's options.
        \else
            \ifnum\lay=\Nnodlen
                \node[layer label] at (\layer,3.6)
                    {Output layer\\$\ell=\layer$\\
                     $\phi^{(\layer)}=\Activation{\layer}$};
                % Restore color: add text=ColorOutput!60!black to these options.
            \else
                \node[layer label] at (\layer,3.6)
                    {Hidden layer\\$\ell=\layer$\\
                     $\phi^{(\layer)}=\Activation{\layer}$};
                % Restore color: add text=ColorHidden!60!black to these options.
            \fi
        \fi
    }

    % These equations describe the computation represented by each neuron.
    % n_(ell-1) counts neurons in the preceding layer, excluding biases.
    % A bias is its own additive parameter, not an ordinary activation.
    % The elementwise convention applies to the activation list used here.
    \node[formula] at (2,-3.55) {
        $z_i^{(\ell)}=
          \displaystyle\sum_{j=1}^{n_{\ell-1}}
             w_{ij}^{(\ell)}a_j^{(\ell-1)}+b_i^{(\ell)},
        \qquad a_i^{(\ell)}=\phi^{(\ell)}\!\left(z_i^{(\ell)}\right)$
    };
    \node[formula] at (2,-4.15) {
        $\mathbf{z}^{(\ell)}=W^{(\ell)}\mathbf{a}^{(\ell-1)}
             +\mathbf{b}^{(\ell)},
        \qquad \mathbf{a}^{(\ell)}=
             \phi^{(\ell)}\!\left(\mathbf{z}^{(\ell)}\right)
             \quad\text{(elementwise)}$
    };
    \node[formula] at (2,-4.65) {
        $\ReLU(u)=\max\{0,u\},\qquad
          \tanh(u)=\dfrac{e^u-e^{-u}}{e^u+e^{-u}},\qquad
          \sigma(u)=\dfrac{1}{1+e^{-u}}$
    };
    % The last line defines the example activation functions above.
    % If you introduce another function, update this explanatory line too.

    % PASS 2: connect each destination i to each preceding source j.
    \foreachitem \N \in \Nnod{
        \edef\lay{\Ncnt}
        \pgfmathtruncatemacro{\prev}{\lay-1}
        % \prev is the source's 1-based LIST position. Numerically it
        % also equals the destination's 0-based mathematical layer \ell.
        \ifnum\lay>1 % Input layer has no incoming weight matrix.
            \foreach \i in {1,...,\N}{
                % Different destinations use different label positions.
                % This staggers the values horizontally to avoid overlap.
                \pgfmathsetmacro{\labelpos}{0.16+0.64*(\i-1)/max(1,\N-1)}
                \foreach \j in {1,...,\Nnod[\prev]}{
                    % Background strokes cannot overprint labels or neurons.
                    \begin{pgfonlayer}{background}
                        \draw[connect arrow] (N\prev-\j) -- (N\lay-\i);
                    \end{pgfonlayer}
                    % -- means a straight segment. pos=0 is the start;
                    % pos=1 is the end. The old pos=0.69 meant 69% along.
                    % A path node writes text at that fractional position.
                    % fill=white masks crossing strokes behind the text.
                    \path (N\prev-\j) --
                        node[weight label,pos=\labelpos]
                            {$\Weight{\prev}{\i}{\j}$}
                        (N\lay-\i);
                    % This lookup replaces the original hard-coded {100}.
                    % Even zero-valued weights retain a labeled edge.
                }
            }
        \fi
    }

    % In the original (N1-1.90), N1-1 names the node and .90 names its
    % boundary anchor at angle 90 degrees (the top), NOT neuron 1.90.
    % above=5 meant place text 5pt above that anchor.
    % Fixed heading coordinates above avoid repeating that anchor syntax.
    % \message{...} writes compiler-log text; ^^J inserts a log newline.
    % It does not put text into the diagram, so log messages were omitted.
    % The original contour and algorithmic packages were unused and removed.
    % References for the explained syntax:
    % https://tikz.dev/tikz-arrows
    % https://tikz.dev/tikz-shapes
    % https://ctan.org/pkg/listofitems
\end{tikzpicture}
\end{document}
